\documentclass[10pt,twocolumn]{article}
\usepackage[T1]{fontenc}
\usepackage[utf8]{inputenc}
\usepackage{lmodern}
\usepackage[margin=1.8cm,top=2.4cm,bottom=2cm,columnsep=0.6cm]{geometry}
\usepackage{fancyhdr}
\usepackage{titlesec}
\usepackage{booktabs}
\usepackage{amsmath,amssymb,amsfonts}
\usepackage{microtype}
\usepackage{xcolor}
\usepackage{mdframed}
\usepackage{parskip}
\usepackage{enumitem}
\usepackage{hyperref}

\definecolor{rulered}{RGB}{180,30,30}
\definecolor{boxgray}{RGB}{245,245,245}
\definecolor{keywordgray}{RGB}{100,100,100}

\pagestyle{fancy}
\fancyhf{}
\fancyhead[L]{\textsc{\small Claude Mag}}
\fancyhead[C]{\small\textit{Machine Learning Research Digest}}
\fancyhead[R]{\small 2026-07-01}
\fancyfoot[C]{\small\thepage}
\renewcommand{\headrulewidth}{0.8pt}

\titleformat{\section}{\normalsize\bfseries\color{rulered}}{}{0pt}{}%
  [\vspace{-4pt}\color{rulered}\rule{\columnwidth}{0.4pt}]
\titlespacing{\section}{0pt}{8pt}{4pt}

\hypersetup{colorlinks=true,urlcolor=rulered,linkcolor=black}

\begin{document}
\setlength{\parindent}{0pt}
\setlength{\parskip}{4pt}

{\small\color{keywordgray}\textbf{Keywords:}
  single-image generation \textbullet{}
  diffusion models \textbullet{}
  patch statistics \textbullet{}
  training-free}\par\vspace{4pt}

{\LARGE\bfseries\raggedright
  Efficient and Training-Free Single-Image Diffusion Models}\par\vspace{4pt}

{\large\itshape\raggedright
  Closed-Form Patch Denoising Replaces Hours of Training with
  Seconds of Inference---CVPR 2026 Highlight}\par\vspace{6pt}

{\small\textbf{By} Haojun Qiu, Kiriakos N.\ Kutulakos \& David B.\ Lindell
  (University of Toronto / Vector Institute)}\par\vspace{2pt}

{\small\color{keywordgray}arXiv:2606.04299 \textbullet{}
  CVPR 2026 Highlight \textbullet{} Submitted June 2026}
\par\vspace{2pt}\noindent{\color{rulered}\rule{\columnwidth}{0.6pt}}\vspace{6pt}

Single-image generative models no longer need to be trained. By treating a
reference image as an empirical patch distribution and computing the score
function analytically in closed form, a new method from the University of
Toronto generates diverse, high-quality samples in seconds---$300\times$ faster
than trained alternatives---while matching or exceeding their quality on
standard benchmarks.

\begin{mdframed}[backgroundcolor=boxgray,linecolor=rulered,linewidth=1pt,
    innerleftmargin=6pt,innerrightmargin=6pt,
    innertopmargin=6pt,innerbottommargin=6pt]
\textbf{\small AT A GLANCE}\par\vspace{4pt}
\begin{description}[leftmargin=0pt,labelsep=4pt,itemsep=2pt,parsep=0pt,
    font=\small\bfseries]
  \item[Problem:] Generating novel images that match the internal patch
    statistics of a single reference image currently requires hours of neural
    network training.
  \item[Why it matters:] Training-based single-image methods are too slow for
    interactive editing, personalised synthesis, or low-resource scenarios;
    a training-free alternative removes the bottleneck entirely.
  \item[Method:] Model the image as a finite empirical distribution over patches
    at multiple scales; compute the diffusion score function in closed form via
    Gaussian kernel regression over that patch set.
  \item[Result:] State-of-the-art SIFID on most textures, $300\times$ speed-up
    over SinDDPM, and support for text-guided stylisation and retargeting without
    any fine-tuning.
  \item[Evidence:] CVPR 2026 Highlight; evaluated against SinDDPM, SinFusion,
    SSGAN, ConSinGAN on standard single-image benchmarks.
\end{description}
\end{mdframed}

\section{The Big Picture}

Single-image generation (SIG) asks: given one photograph or texture, synthesise
new images that share its internal visual statistics. Applications span
texture synthesis, artistic style transfer, image retargeting, and
data augmentation for rare visual classes. The task is intrinsically
self-supervised---the only training signal is the image itself.

The dominant recent approach trains a small diffusion model on the patches
extracted from a single image (SinDDPM, SinFusion). While effective, even
these small networks require hours of GPU time per reference image---a steep
price when the target is a single photo. The underlying cost is neural
network training, which is needed because the score function is otherwise
intractable for a general data distribution.

This paper observes that for a single image the score function \emph{is}
tractable: the image contains only finitely many patches, so the score at any
noisy point is exactly the posterior mean over that finite set---a
Nadaraya-Watson kernel regression with a Gaussian kernel, computable in
closed form. Embedding this closed-form denoiser inside a standard diffusion
sampling loop yields a complete, training-free generative model that runs in
seconds on a commodity GPU.

\section{Why Existing Approaches Fall Short}

Neural single-image diffusion models (SinDDPM, SinFusion) require
training a U-Net on the reference image from scratch. Training takes
1--8 hours depending on image resolution, making interactive use
impractical. Moreover, these networks are \emph{not} reusable: each new
reference image requires fresh training.

Score function engineering approaches (e.g., SinGAN) use a GAN
pyramid but suffer from training instability and do not support
the flexible text-guided generation that diffusion models enable.

The fundamental gap is the absence of a principled, efficient score
estimator that does not require parametric training.

\section{Core Mechanism}

For a single image $I$, extract patches $\mathcal{P} = \{x_1,\ldots,x_N\}$
at multiple scales. For a noisy patch $y = x^* + \sigma\varepsilon$,
$\varepsilon\sim\mathcal{N}(0,I)$, the optimal minimum-MSE denoiser is the
posterior mean:
\[
D_\sigma(y) = \mathbb{E}[x \mid y]
= \frac{\sum_{i=1}^{N} x_i\,e^{-\|y-x_i\|^2/2\sigma^2}}
       {\sum_{i=1}^{N} e^{-\|y-x_i\|^2/2\sigma^2}}.
\]
This is Gaussian kernel regression over the patch set---no neural network
required. The score function follows from Tweedie's formula:
\[
\nabla_y \log p_\sigma(y) = \frac{D_\sigma(y) - y}{\sigma^2}.
\]

\section{Technical Formulation}

\textbf{Patch dataset.}
For resolution $h\times w$ and patch size $p\times p$, the patch set
$\mathcal{P}^\ell$ at scale $\ell$ has $N^\ell = (h_\ell - p + 1)(w_\ell - p + 1)$
elements, each of dimension $d = p^2 \cdot C$ (with $C$ channels).

\textbf{Optimal denoiser.}
Under the empirical prior $p(x) = \frac{1}{N}\sum_{i=1}^N \delta(x - x_i)$,
the posterior given noisy observation $y$ is:
\[
p(x \mid y) \propto \mathcal{N}(y; x, \sigma^2 I)\cdot p(x)
= \sum_{i=1}^N w_i(y)\,\delta(x-x_i),
\]
\[
w_i(y,\sigma) \propto \exp\!\left(-\frac{\|y - x_i\|^2}{2\sigma^2}\right).
\]
The posterior mean $D_\sigma(y) = \sum_i w_i x_i$ is the closed-form denoiser.

\textbf{Computational cost.}
Each denoising step costs $O(N d)$ inner products. For $p=8$, $C=3$
and a $512\times512$ image, $N\approx2.5\times10^5$ and $d=192$,
giving $\approx5\times10^7$ operations---substantially cheaper than a
forward pass of even a small U-Net.

\textbf{Multi-scale generation.}
Coarse-to-fine: generate at the lowest resolution first using
$\mathcal{P}^1$, upsample, and condition the next-scale denoiser on the
upsampled result via score guidance:
\[
\tilde\nabla_y \log p_\sigma^\ell(y) =
\nabla_y\log p_\sigma^\ell(y) + \omega\,\nabla_y\log p_\sigma^{\ell-1}(y_\downarrow),
\]
where $y_\downarrow$ is $y$ downsampled to scale $\ell{-}1$ and $\omega$
is a guidance weight.

\section{Learning or Inference Procedure}

No training is performed. At inference:
\begin{enumerate}[itemsep=2pt,parsep=0pt]
  \item Extract patch sets $\{\mathcal{P}^\ell\}_{\ell=1}^{L}$ from $I$.
  \item Initialise $x_T^1 \sim \mathcal{N}(0,I)$ at the coarsest scale.
  \item Run DDPM reverse process using $D_\sigma$ as the denoiser (instead
    of a neural network) for $T$ steps.
  \item Upsample, condition, and repeat for finer scales.
  \item Output $x_0^L$ as the generated image.
\end{enumerate}
Total wall-clock time: $<10$ seconds on a single GPU for a $256\times256$
output, vs.\ $>3$ hours for SinDDPM.

\section{What the Guarantee Says}

The paper proves that $D_\sigma$ is the \emph{Bayes-optimal} denoiser under
the empirical patch prior: no estimator achieves lower MSE for that prior.
As $N\to\infty$ (more patches, i.e., larger image), the empirical prior
concentrates on the true patch distribution and $D_\sigma$ converges to
the score function of the true distribution---so the method is
asymptotically exact.

\section{Experimental Findings}

On the Oxford Textures and USC-SIPI image datasets:
\begin{itemize}[itemsep=2pt,parsep=0pt]
  \item Lower SIFID (Single-Image FID) than SinDDPM on
    8 of 10 tested textures.
  \item $300\times$ faster at inference (seconds vs.\ hours).
  \item Competitive FID on random sample diversity.
  \item Text-guided stylisation (using CLIP guidance on the denoiser score)
    produces coherent results with no fine-tuning.
  \item Image retargeting (change aspect ratio while preserving texture)
    outperforms SinGAN in perceptual quality.
\end{itemize}

\section{Ablations and Interpretation}

\textbf{Patch size $p$:} $p=8$ gives the best SIFID/diversity trade-off;
larger patches lose texture variety (fewer distinct patches), smaller
patches lose global coherence.

\textbf{Number of scales $L$:} $L=3$ is sufficient; $L=1$ generates
texturally correct but spatially incoherent images.

\textbf{Guidance weight $\omega$:} Higher $\omega$ produces images more
faithful to the reference layout; lower $\omega$ increases diversity.

\section*{Reference}

Haojun Qiu, Kiriakos N.\ Kutulakos, David B.\ Lindell.
\textbf{Efficient and Training-Free Single-Image Diffusion Models}.
CVPR 2026. arXiv:2606.04299.
\url{https://arxiv.org/abs/2606.04299}

\end{document}
